\documentclass[11pt,a4paper]{amsart}

\usepackage{hyperref}
\usepackage{fullpage}
\usepackage{etoolbox}
\usepackage[british]{babel}


%%\usepackage{titlesec}
%%\titleformat*{\section}{\center\sffamily}


\patchcmd{\section}{\normalfont}{\normalfont\Large}{}{}
%\patchcmd{\section}{\scshape}{\bfseries}{}{}
%\makeatletter
%\renewcommand{\@secnumfont}{\normalfont}
%\makeatother

\renewcommand{\labelenumi}{\theenumi.}
%\renewcommand{\thesubsection}{\Alph{subsection}}

\linespread{1.2}

\setlength{\parindent}{0pt}
\setlength{\parskip}{5pt}


\begin{document}
\begin{center} 

{\bf \sffamily \Large Guideline for Tutors} \\ 
\vspace{3mm} 
School of Mathematical Sciences \\ 
Lancaster University \\
Last updated on \today \\
\end {center}

\noindent\makebox[\linewidth]{\rule{\textwidth}{0.4pt}}
\vspace{2mm}

Thank you for offering to be a tutor (aka GTA = Graduate Teaching Assistant) this year. The department prides itself on the quality of its teaching and greatly values the enthusiastic support of its tutors. We hope that your involvement in this activity will be rewarding and provide you with valuable career experience. This document contains some guidance, and outlines both the duties of tutors as well as the support that they can expect from the department in return.

\section{Scope of this document}

This document refers to tutors helping with workshops within the undergraduate programme in Mathematics. It is primarily intended for Graduate Teaching Assistants (GTAs); however it also contains useful information for other staff. Most workshops for modules with the MATH prefix are covered. Some modules such as MATH240T have a less standard structure, and modules for other departments may have different or additional guidelines.


\section{Workshops}

Most undergraduate modules have workshops (weekly or fortnightly). Typically workshops have a dual purpose: 
\begin{itemize}
 \item to help students work through the ``workshop questions'', and
 \item to give an opportunity for one-on-one feedback on their assessed coursework (which may include online quiz questions). 
\end{itemize}
Unless instructed otherwise by your lecturer, the following {\bf advice} will apply to your workshops.

\begin{enumerate}
\item You can find all course materials such as lecture notes, workshop questions, etc. on the Moodle page of the corresponding module. 
\item {\bf Familiarise yourself} with the lecture notes and workshop questions (and solutions) beforehand. If you are unsure of anything, clarify it with the lecturer in advance (this is infinitely better than hiding your confusion!). 
\item Turn up at the workshop room {\bf on time}. Bring whiteboard pens. 
You can find them in the cupboard opposite the School Office in the social area on C floor.
\item Use the projector in the room to display the workshop questions during the workshop. 
If the projection device is not suitable then you can also print the questions for them.
\item Attendance monitoring. Students have to register their attendance via the iLancaster app. 
	  Occasionally, it is not working and students ask you to register them manually. 
	  You can do this either via the iLancaster app or online via
	  \url{https://www.lancaster.ac.uk/itpi/ilancaster}
  	 and then selecting `teaching class attendance'. 
  	 You can do that also after the session and also remove students.
  	 However, do not spend too much time on this, since it is not your job to do attendance monitoring. Feel free to direct such students elsewhere, such as the Maths Student Office.
\item Typically we would expect you to spend at most about 15 minutes handing back coursework, providing individual feedback, and going through common mistakes on the whiteboard. The majority of the workshop should be spent on the workshop exercises. 
\item When providing feedback on assessed coursework, even though model solutions will be provided by the lecturer on Moodle, be aware that some students have difficulty in following solutions that are presented without verbal explanation, or in seeing how incorrect solutions are deficient.
\item {\bf Actively circulate}, ensuring that you make contact with each student at least once or twice (beyond giving feedback on assessed work). Some students are reluctant to raise their hands, and can waste an entire workshop stuck on a single question and not ask. Others may assume that their solution is satisfactory when it is not.
\item {\bf Coach don't criticise}: Be encouraging and avoid being patronising or arrogant when answering questions or correcting mistakes.
\item {\bf Give suitable hints}. When a student is stuck, a skilled teacher will give the minimal hint needed that will result in the student answering (most of) the question for themselves. A hint is not necessarily the next step, but instead could be advice or a question which prompts them to figure out the next step for themselves. [The amount of help heavily depends on the student, so this is not an easy skill to master.]
\item There have been very few incidents involving misconduct of students in teaching sessions. If a student is acting in a manner which is unacceptable or threatens to become unacceptable, then you should ask the student to leave. Conversely, you should not do anything to provoke trouble with students.
\end{enumerate}
 \newpage
\section{Marking coursework}

Coursework is typically intended to:
\begin{itemize}
\item Help students to learn the course material by tackling specific problems.
\item Teach students how to write mathematics. 
\item Give feedback to students on their own understanding.
\item Give feedback to the lecturer on student understanding.
\item Inform workshop tutors of specific areas of weakness that a student may need help with.
\item Contribute to the overall assessment of the module.
\end{itemize}

The following advice should be followed, unless the marking scheme indicates otherwise:
\begin{enumerate}
\item Most coursework is written on paper. 
The students have to hand them in the pigeonholes that are currently on A floor in the Charles Carter Building (but this location is subject to review, so look out for an annoucements on their relocation).
The deadlines can be found on Moodle.
\item You should mark them and bring them to the next workshop. If the timing is too tight, please contact the lecturer. 	  
\item {\bf Don't dwell} on partial solutions or spend time hunting for minor errors; incomplete or incorrect solutions can be given partial credit for demonstrating some relevant knowledge or understanding. If a student makes an error near the start of the question, quickly evaluate the amount of understanding they have demonstrated, give it a mark, and move on.
\item To encourage good writing, you should {\bf point out notational and other minor errors} even when you do not penalize them, and conversely you should indicate when a solution is well presented.
\item  It is the students' responsibility to produce work which is {\bf legible}. However, some students have learning needs which may have an impact on their writing and spelling. You should therefore be somewhat flexible when marking work with imperfect handwriting. If the handwriting is hard to decipher, then you should indicate any passages which you do not understand. Any student who regularly produces such work should be advised to improve their presentation, or to submit typeset work.
\item If a coursework is not readable at all, you can reject it. 
      However, before you do that check with the lecturer as the student might has an ILSP.
\item Please add {\bf helpful comments} to students' scripts. However, there is no need to repeat substantial amounts of information that is contained in specimen solutions since these will be made available to the students. Do not be sarcastic or arrogant when writing comments on poor quality work.
\item  If a {\bf mistake is common} to many students, inform the lecturer and consider discussing this point with the whole workshop room.
 \item  If you are ever {\bf unsure what marks to give} for any particular question then you must consult the module lecturer. Sometimes, groups of tutors find it helpful to discuss their marking collaboratively.
\item Logically correct and complete solutions should gain credit, {\bf even if they differ} from the specimen solution. Some students attempt solutions using methods that are totally inconsistent with the marking scheme. In such cases try to award the overall marks appropriately, but if you are not immediately sure what is suitable, please bring the script to the lecturer for advice.
\item For most modules, you have to enter the marks also on Moodle for in a coursework submission point. 
If you are unsure how to do that, please ask for advice.
\item {\bf Late work}: When coursework is submitted after the deadline and without an agreed extension, it should lose 10 per cent of the maximum mark available. Tutors should directly apply this penalty to the piece of work, and also indicate on the mark sheet that a late penalty has been applied. Where the coursework is submitted after the work has already been marked and returned to other students in the course, it should be awarded zero.
\item {\bf Lost work}: It is the student's responsibility to ensure that you receive their work. If a student claims they have submitted work to your pigeonhole, but you didn't received it, then it is not your responsibility to hunt it down. You can direct further student queries about this to the module convenor. In particular, please do not send a mass email asking colleagues to check their pigeonholes.
(NB. Some modules have online homework submission while others have paper submission. Module lecturers should brief you in each case.)
\item {\bf Plagiarism}: This occurs whenever a student appropriates the writings of others (including AI), and dishonestly presents these as if they were the student's own; usually this happens when they fail to distinguish between common knowledge and knowledge that should be attributed to others. Students are welcome to discuss work with other students, but all pieces of work
that they submit must be their own words and thoughts. If you think that work has been produced improperly, you should contact the lecturer in the first instance, who may then contact the Academic Integrity Officer (in 2026/27, that's Daniel Elton). Ensure you retain the relevant evidence. Similarities between coursework of different students do not necessarily indicate improper collusion; many students have been taught to write mathematics in similar styles.
\end{enumerate}

\newpage

\section{The Department's obligation to you}

\begin{enumerate}
\item  The department will inform you of your teaching duties at the beginning of the term and will issue you with the necessary forms and documentation for your employment.
\item The department will process your claim form promptly so that you receive payment for your teaching in a timely way.
\item The department will provide you with appropriate {\bf training} for the teaching activities that you are involved in.
This will include the {\bf Introduction to Teaching and Lancaster (ITL)}:
\begin{center} \url{https://www.lancaster.ac.uk/curriculum-and-education-development-academy/workshops/introduction-to-teaching-for-lancaster-itl/} \end{center} \par
It is likely that the department will (probably in Michaelmas term) run a {\bf Marking Calibration exercise} to help to improve the consistency and speed of your marking, or additional {\bf Teaching Enhancement workshops} for GTAs. These are paid training opportunities.
\item {\bf Homework pigeonholes}, into which students submit their coursework for marking, are organised by workshop groups. The department will provide you with keys to access the pigeonholes for any groups you are teaching. 
(NB. Some modules have online homework submission.)
\item  You will be given access to the Moodle pages for any modules you teach. This is usually the best place to access course materials (lecture notes, exercise sheets, solutions, marking schemes, deadline), but you may also hear from the module lecture.
\item The module lecturer will be available for questions or clarification about anything related to the module, including your role in it, content questions, or rationale behind the assessment methods.
\item The module lecturer will {\bf moderate} your marking when requested. This can be if you are unclear how to apply the marking scheme in a specific case or if a student queries the mark that they have been allocated. The module lecturer may also ask to inspect your marking in order to maintain consistency and quality, and to provide mentoring.
\item The module lecturer will endeavour to ensure marking can be completed within the {\bf allocated marking time}. For example, typically tutors are paid each week for: 1 hour in workshop attendance and 2 hours on (preparation + marking) per workshop group. Time spent may be somewhat uneven from week to week, but over a whole module it should balance out. 2 hours is the marking allocation for Year 1 midterm tests. In the case of workshops for which there is no marking, the expectation should be for one hour of preparation for the session, plus the number of hours of the session. Although it is always possible to spend more time on an activity in order to deliver it to a higher standard, you should aim to provide the best possible service to students {\bf within the allocated time}. 
\item If you find that it is impossible to carry out the work within the given time, you should discuss the problem with the lecturer, who should find ways of reducing the workload. If the problem persists, please contact the GTA Coordinator (in 2026/27, that's David Pauksztello), or your GTA representative.
\item You are under no obligation to give advice outside the workshop or study material beyond the confines of the workshop questions; {\bf redirect such queries to the lecturer}, who has regular office hours. The department does not regard graduate teaching assistants as having any pastoral role towards undergraduates.
\item The module lecturer will respond to any feedback that you provide to them concerning the aptness of the coursework, the quality of the mark scheme and the structure of the workshops. Lecturers should discuss and review the amount of work required by tutors at the end of each module with a view to making revisions for the following year.
\item The procedure for recording marks in each module will be explained by the lecturer or the department Teaching Office. (For online submission, this will typically happen automatically, once you have recorded marks on moodle.)
\end{enumerate}

\section{Your obligations}

\begin{enumerate}
\item You should provide the Departmental Officer (Sarah Milne) with any information that is requested concerning your employment.
%\item You must attend the training sessions required by the department. In particular, the {\bf Introduction to Teaching and Lancaster (ITL)}:
%\begin{center} \url{http://www.lancaster.ac.uk/oed/academic-and-research-staff/introduction-to-teaching-at-lancaster/} \end{center} \par
%Attendance will be recorded and GTAs who have not attended this or an equivalent training session may be prevented from tutoring. Also, in Michaelmas term there is a {\bf Marking Calibration exercise} in the department, which each GTA must attend at least once every two years. 
\item You should follow any instructions issued to you by the module lecturer about how they would like you to participate in their module.
\item When you sign up for teaching, you should be committed to teaching your full allocated portion of the module (modulo sickness or one-off absences such as conference attendance). We do not offer the option of switching out of a module after you have agreed to take a workshop except in exceptional circumstances (e.g. longer-termn sickness).
\item You should ensure that you are sufficiently {\bf familiar with the course materials} so that you can competently mark the coursework, and answer questions that students are likely to put to you during workshop sessions.
\item You must carefully follow the lecturer's {\bf mark scheme} and use the specimen solutions. If anything is unclear, you are expected to ask the lecturer for clarification.
\item You will {\bf record all your marks} in the required format {\bf for each assessment}. If you fail to submit the marks without providing a good reason, the module lecturer and your line manager (typically, PhD supervisor) will be informed. If you fail to return the completed mark sheet within two weeks of the end of the module, the Head of Department will be informed. It is your responsibility to ensure that mark sheets are complete and correct.
\item You must {\bf return the marked work} to the students on time, typically this is during the workshop that immediately follows the submission deadline. Plan ahead to ensure that is achieved. If you are handing over your tutor group to another tutor at the end of a module or because of a planned absence, you must ensure that they receive the marked work before the workshop in which it is to be returned to the student. If you are not expecting to see your tutor group again, you should mark the work promptly (for regular weekly assignments, no more than 4 working days) and hand it to the module convenor, who will be able to return it to the students. Return uncollected work to the module convenor at the end of the module.
\item You should know when and where your workshops are. (These should appear on your online Timetable - please enquire with the Student Office if you think your scheduled workshops are not the same as those you agreed to teach.)
\item You must return your {\bf claim sheet} to the department officer promptly. Any claims larger than the standard amount allocated for your the module must be approved by the Director of Undergraduate Teaching.
\item If you are unable to mark an assignment or attend a workshop session (unless it is due to illness, see below), you must {\bf nominate a substitute}, provided that:
\begin{itemize}
\item You confirm these arrangements with the module lecturer and the Student Office;
\item The substitute is a current (trained) tutor or has taught on the course before; 
\item The substitute tutor is provided with sufficient teaching materials (notes, specimen solutions, access to students’ submitted coursework etc.) to cover the assignment;
\item The substitute tutor understands whether they are only substituting for the workshop, or whether you need them to collect and carry out your marking for you as well (in which case they will need access to the workshop pigeonhole).
\end{itemize}
In most cases, {\bf it is your responsibility to find a substitute}, and essential that you ensure that adequate cover is being provided in your absence.
\item If you are unable to attend a workshop or mark an assignment due to {\bf illness}, then please inform the Student Office a.s.a.p.~ so that we can find appropriate cover. You are not expected to find a substitute yourself, in this case.
\end{enumerate}



If you have any queries or concerns relating to tutoring, please feel free to contact me.

\begin{flushright} 
Mark MacDonald \\ 
Director of Undergraduate Teaching \\
Mathematical Sciences \\

\bigskip

David Pauksztello \\
GTA Coordinator \\
Mathematical Sciences
\end{flushright}

\end{document}